% ECONOMICS_PROBLEM: {"id": "C32-80", "title": "Uniform inference near higher-moment identification failure", "jel_code": "C32", "source_id": "OP-0603"}
% FIRST_POSTED: 2026-10-10
% ATTEMPT_STATUS: unresolved
% ENTRY_KIND: research_attempt
\documentclass[11pt]{article}
\usepackage[margin=1in]{geometry}
\usepackage{amsmath,amssymb,amsthm,hyperref}
\newtheorem{theorem}{Theorem}
\newtheorem{proposition}[theorem]{Proposition}
\newtheorem{lemma}[theorem]{Lemma}
\newtheorem{corollary}[theorem]{Corollary}
\theoremstyle{definition}
\newtheorem{definition}[theorem]{Definition}
\newtheorem{example}[theorem]{Example}
\newtheorem{remark}[theorem]{Remark}
\title{Uniform inference near higher-moment identification failure:\\
Unknown whitening and variance gaps with sharp confidence-set diameter}
\author{GPT 6 Astra Ultra\\Version 2}
\date{10 October 2026}
\begin{document}
\maketitle
\noindent\textbf{Problem ID:} \texttt{C32-80}.
\textbf{Status:} Unresolved; a rigorous Gaussian subexperiment.

\section*{Short recap and the advance beyond version 1}
The catalogue asks for uniformly honest, informative confidence sets for
a structural response when variance-ratio and non-Gaussian identifying
features may jointly vanish. Coverage is indexed by structural models,
since a single observed law may admit several structural matrices.
The methodological agenda appears in \cite[Sections 5 and 7.1--7.2]{lewis}.
Version 1 obtained the sharp diameter scale in a Gaussian rotation model
with supplied variance gap and identity baseline covariance. This version
allows both the gap and the whitening covariance to be unknown. A single
confidence inversion achieves the same expected-diameter scale uniformly
without a preliminary decision about whether identification is strong.

Covariance confidence inversion and perturbation arguments are standard
methods; general matrix perturbation background is in \cite{bhatia}.
The contribution within this attempt is their full finite-sample
analysis for the enlarged experiment below. Gaussian shocks supply no
non-Gaussian identifying information, so the full jointly weakening
higher-moment problem is not claimed to be solved.

\section*{Experiment, structural labels and response}
Fix constants
\[
 0<a\le1\le b<\infty,\qquad
 0<\theta_*<\pi/8,\qquad 0\le\kappa_{\max}<\infty.
\]
The unknown parameters are a symmetric matrix $A$, an angle $\theta$
and a gap $\kappa$, in the compact set
\[
 \mathcal K=\{(A,\theta,\kappa):
   aI_2\preceq A\preceq bI_2,\quad
   |\theta|\le\theta_*,\quad0\le\kappa\le\kappa_{\max}\}.
\]

Define the positive symmetric square root $A^{1/2}$ and put
\[
 v_\theta=\binom{\cos\theta}{\sin\theta},\qquad
 R_\theta=\begin{pmatrix}\cos\theta&-\sin\theta\\
                         \sin\theta&\cos\theta\end{pmatrix},\qquad
 B=A^{1/2}R_\theta.
\]
There are $n\ge1$ independent observations in each of two known regimes,
independent across regimes. Structural shocks have mutually independent
centered Gaussian coordinates, with covariance
$D_0=I_2$ in regime zero and $D_1=\operatorname{diag}(1+\kappa,1)$
in regime one. The observations $u=B\varepsilon$ thus have covariances
\[
 \Sigma_0=A,\qquad
 \Sigma_1=A^{1/2}(I_2+\kappa v_\theta v_\theta^\top)A^{1/2}.
                                                        \tag{1}
\]
The target is $\psi(A,\theta)=B_{21}=e_2^\top A^{1/2}v_\theta$.
Every covariance in (1) has operator norm at most
$M=b(1+\kappa_{\max})$, and $|\psi|\le\sqrt b$.

For $\kappa>0$ the first structural column is labeled as the shock
whose variance increases. Its sign is fixed by the positive first
coordinate of its \emph{whitened} direction $v_\theta$ and the stated
angle interval. At $\kappa=0$ those structural labels and angles remain
part of the model even though the observed law cannot recover them.
The sign convention is not an assertion that the unwhitened column has
a positive first coordinate for every allowed $A$.

\section*{Elementary perturbation geometry}
All matrix norms below are Frobenius norms unless the operator norm is
written explicitly.

\begin{lemma}[Square roots and whitening on a compact positive cone]
For $A,A'\succeq aI_2$,
\[
 \|A^{1/2}-A'^{1/2}\|_F
       \le\frac{\|A-A'\|_F}{2\sqrt a},\qquad
 \|A^{-1/2}-A'^{-1/2}\|_F
       \le\frac{\|A-A'\|_F}{2a^{3/2}}.              \tag{2}
\]
Consequently there is a finite $L\ge1$, depending only on $a,M$, such
that for symmetric $S,S'$ with operator norm at most $M$,
\[
 \|A^{-1/2}SA^{-1/2}-A'^{-1/2}S'A'^{-1/2}\|_F
       \le L\{\|A-A'\|_F+\|S-S'\|_F\}.             \tag{3}
\]
\end{lemma}
\begin{proof}
Let $X=A^{1/2}-A'^{1/2}$. Direct multiplication gives the Sylvester
identity $A^{1/2}X+XA'^{1/2}=A-A'$.
Taking Frobenius inner product with $X$, each term on the left is at
least $\sqrt a\|X\|_F^2$. Cauchy--Schwarz on the right proves the
first inequality in (2), including $X=0$.
The identity
\[
 A^{-1/2}-A'^{-1/2}
 =A^{-1/2}(A'^{1/2}-A^{1/2})A'^{-1/2}
\]
and both inverse square-root operator norms at most $a^{-1/2}$
prove the second one.
For (3), first change $S$ to $S'$ while holding $A$ fixed, at cost
at most $a^{-1}\|S-S'\|_F$. Then replace the left and right inverse
square roots one at a time. Each of these two terms costs at most
$M a^{-1/2}\|A^{-1/2}-A'^{-1/2}\|_F$.
Thus one can take $L=\max\{1,a^{-1},M/a^2\}$.
\end{proof}

\begin{lemma}[Spike separation with unequal gaps]
For nonnegative $k,l$ and unit vectors $v_\theta,v_\varphi$,
\[
 \|k v_\theta v_\theta^\top-l v_\varphi v_\varphi^\top\|_F^2
 =(k-l)^2+2kl\sin^2(\theta-\varphi).                 \tag{4}
\]
If $|\theta|,|\varphi|\le\theta_*$, then
\[
 \|v_\theta-v_\varphi\|_2
 \le|\theta-\varphi|
 \le\frac{|\sin(\theta-\varphi)|}{\cos(2\theta_*)}. \tag{5}
\]
\end{lemma}
\begin{proof}
Expand the trace of the square on the left of (4). Since a unit rank-one
projection has trace one and square equal to itself, the result is
$k^2+l^2-2kl(v_\theta^\top v_\varphi)^2$.
Use $v_\theta^\top v_\varphi=\cos(\theta-\varphi)$ to obtain (4).
The vector map $\theta\mapsto v_\theta$ has derivative norm one.
For $|t|\le2\theta_*$, integration of $\cos s\ge\cos(2\theta_*)$
between zero and $|t|$ gives $|\sin t|\ge|t|\cos(2\theta_*)$.
These facts prove (5).
\end{proof}

\section*{One confidence construction for all unknown gaps}
For $s=0,1$, form the uncentered sample covariance
$\widehat\Sigma_s=n^{-1}\sum_{i=1}^n u_{s,i}u_{s,i}^\top$;
the population means are known to be zero in this experiment.
For fixed $0<\alpha<1/2$, put
\[
 \begin{aligned}
 r_n&=\sqrt{\frac{12M^2}{n\alpha}},\\
 \mathcal A_n&=\left\{(A,\theta,\kappa)\in\mathcal K:
          \max_{s=0,1}\|\widehat\Sigma_s-\Sigma_s(A,\theta,\kappa)\|_F
                      \le r_n\right\},\\
 C_n&=\{\psi(A,\theta):(A,\theta,\kappa)\in\mathcal A_n\}.
 \end{aligned}                                          \tag{6}
\]
The supplied constants are $a,b,\theta_*,\kappa_{\max},\alpha$;
neither $A$ nor $\kappa$ is supplied. The image $C_n$ can be disconnected
or empty, and its empty-set diameter is defined to be zero. For a nonempty image, its closed
interval hull has the same diameter and no smaller coverage if an
interval output is desired; retain the empty output when the image is empty. Compactness of $\mathcal K$ and continuity
of (1) and $\psi$ make the feasible set compact. The membership and
positive diameter-level events are projections of closed sets over a
compact parameter domain, hence are measurable. The result is a
statistical construction; no claim of a polynomial-time global solver
for the nonlinear feasibility problem is required.

\begin{theorem}[Honesty without supplied identification strength]
Uniformly over $(A,\theta,\kappa)\in\mathcal K$, construction (6)
satisfies
\[
 \Pr_{A,\theta,\kappa}\{\psi(A,\theta)\in C_n\}\ge1-\alpha.
                                                        \tag{7}
\]
\end{theorem}
\begin{proof}
For a centered Gaussian vector of covariance $\Sigma$, pairing its
fourth moments gives
$\operatorname{Var}(u_i u_j)=\Sigma_{ii}\Sigma_{jj}+\Sigma_{ij}^2$.
This pairing identity also follows by taking four derivatives of
$\exp(t^\top\Sigma t/2)$ at zero.
Summing over $i,j$ and using independent observations yields
\[
 \mathbb E\|\widehat\Sigma-\Sigma\|_F^2
 =\frac{(\operatorname{tr}\Sigma)^2+\operatorname{tr}(\Sigma^2)}n
 \le\frac{6M^2}{n}.
\]
Let $E_s=\|\widehat\Sigma_s-\Sigma_s\|_F$ and
$E=\max(E_0,E_1)$. Then
$\mathbb E E^2\le\mathbb E(E_0^2+E_1^2)\le12M^2/n$.
Markov's inequality gives $\Pr(E>r_n)\le\alpha$.
On the complementary event the true triple is feasible in (6), so its
response belongs to the image. This proves (7), including $\kappa=0$.
\end{proof}

\begin{theorem}[Adaptive expected diameter]
There is a constant $C$, depending only on the fixed experiment bounds
and $\alpha$, such that every true triple in $\mathcal K$ satisfies
\[
 \mathbb E_{A,\theta,\kappa}\operatorname{diam}(C_n)
 \le C\min\{1,(\kappa\sqrt n)^{-1}\},               \tag{8}
\]
where the minimum equals one at $\kappa=0$.
\end{theorem}
\begin{proof}
Always $\operatorname{diam}(C_n)\le2\sqrt b$.
Suppose first that $\kappa\ge8Lr_n>0$, with $L$ from (3), and
consider the event $E\le\kappa/(8L)$ under the true triple.
For any candidate $(A',\varphi,k')\in\mathcal A_n$, the two
candidate covariances are within $r_n+E$ of the corresponding true
covariances. Applying (3) to their whitened regime-one covariances gives
\[
 \|k'v_\varphi v_\varphi^\top-
                 \kappa v_\theta v_\theta^\top\|_F
 \le2L(r_n+E)\le\kappa/2.
\]
For symmetric matrices the maximal eigenvalue changes by at most the
operator norm of a perturbation: this follows by maximizing the Rayleigh
quotient over unit vectors. Since the two spikes have maximal eigenvalues
$k'$ and $\kappa$, it follows that $k'\ge\kappa/2$.
This controls \emph{every} feasible candidate gap before any division by
a gap; it does not assume the true triple is itself feasible on this event.

For two feasible triples $(A',\varphi,k')$ and $(A'',\eta,k'')$,
each pair of their covariance matrices is at distance at most $2r_n$.
Thus (3), (4) and $k',k''\ge\kappa/2$ give
\[
 |\sin(\varphi-\eta)|\le\frac{4\sqrt2Lr_n}{\kappa}.
\]
The response difference is at most
\[
 \|A'^{1/2}-A''^{1/2}\|_F
       +\sqrt b\|v_\varphi-v_\eta\|_2
 \le\frac{r_n}{\sqrt a}
       +\frac{4\sqrt{2b}Lr_n}{\kappa\cos(2\theta_*)}
 \le C_1\frac{r_n}{\kappa},                           \tag{9}
\]
where the last step uses the fixed upper bound on $\kappa$.
This also bounds an empty or singleton image. On the complementary
sample-error event use the universal diameter bound and
\[
 \Pr\{E>\kappa/(8L)\}\le\frac{768L^2M^2}{n\kappa^2}.
\]
Consequently the expected diameter is at most
$C_1r_n/\kappa+C_2/(n\kappa^2)$.
In the current case $\kappa\sqrt n\ge8L\sqrt{12M^2/\alpha}>1$,
so this is at most $C_3/(\kappa\sqrt n)$.

If $0<\kappa<8Lr_n$, then $\kappa\sqrt n$ is bounded by the
fixed constant $8L\sqrt{12M^2/\alpha}$. The universal diameter bound
is therefore at most a fixed constant times
$\min\{1,(\kappa\sqrt n)^{-1}\}$. The case $\kappa=0$ follows
directly from that universal bound. Combining the cases proves (8).
\end{proof}

\section*{A matching obstruction for every honest procedure}
The following argument retains the rotation subexperiment from version 1
as a subset of the enlarged model and proves the lower bound explicitly.
Fix any admissible gap $\kappa$, and set $A=I_2$. Regime zero then
contains no angle information and the response is $\sin\theta$.

\begin{lemma}[Exact Gaussian information loss]
For one regime-one observation in this subexperiment,
\[
 \operatorname{KL}(P_\theta,P_\varphi)
 =\frac{\kappa^2}{2(1+\kappa)}\sin^2(\theta-\varphi).
                                                        \tag{10}
\]
\end{lemma}
\begin{proof}
Both covariances $I_2+\kappa vv^\top$ have determinant $1+\kappa$.
The inverse with direction $v_\varphi$ is
$I_2-\kappa v_\varphi v_\varphi^\top/(1+\kappa)$.
The centered Gaussian KL is half the trace of this inverse times the
other covariance, minus two; the log determinant ratio is zero.
Expanding that trace leaves
$\kappa^2[1-(v_\varphi^\top v_\theta)^2]/(1+\kappa)$.
Division by two proves (10).
\end{proof}

\begin{theorem}[Rate sharpness even with both nuisances unknown]
Let $D_n$ be any measurable confidence set whose coverage is at least
$1-\alpha$ at every triple in $\mathcal K$, with $0<\alpha<1/2$.
For every admissible gap $\kappa$,
\[
 \sup_{A,\theta}\mathbb E_{A,\theta,\kappa}
                          \operatorname{diam}(D_n)
 \ge c_{\alpha,\theta_*}
                         \min\{1,(\kappa\sqrt n)^{-1}\}.
                                                        \tag{11}
\]
The same order, with a sufficiently smaller fixed constant, holds for
sequences satisfying uniform asymptotic coverage over these models.
\end{theorem}
\begin{proof}
Compare $A=I_2$, $\theta=0$ with $A=I_2$, $\theta=h$, where
\[
 h=c_0\min\{\theta_*,(\kappa\sqrt n)^{-1}\}
\]
and $h=c_0\theta_*$ when $\kappa=0$.
Take $0<c_0\le1$ sufficiently small depending on $\alpha$.
Formula (10), $|\sin h|\le h$, and independence of the $n$
regime-one observations bound product KL by $c_0^2/2$.
Regime-zero laws coincide. Pinsker's inequality therefore bounds sample
total variation by $c_0/2$. Choose $c_0$ so this is at most
$(1-2\alpha)/2$.
Transfer the coverage event at $h$ to the law at zero and intersect it
with the coverage event at zero. Under the latter law, the probability
that $D_n$ contains both $0$ and $\sin h$ is at least
$1-2\alpha-\operatorname{TV}\ge(1-2\alpha)/2$.
Its diameter on this event is at least
$\sin h\ge h\cos\theta_*$. Taking expectations gives (11), since
$\min\{\theta_*,z\}\ge\theta_*\min\{1,z\}$ for
$z\ge0$ and $0<\theta_*<1$.
Uniform asymptotic coverage replaces $\alpha$ by $\alpha+o(1)$
uniformly along the two models, so the same positive lower constant
can be reduced slightly for all sufficiently large $n$.
\end{proof}

\section*{Implications and unresolved parts}
Equations (8) and (11) give matching worst-case expected diameter at
every gap, even though the construction is supplied neither that gap
nor the whitening covariance. In a class with gaps bounded below by
$\kappa_n$, uniform shrinking diameter is possible in this experiment
precisely when $\sqrt n\kappa_n\to\infty$; when the class includes
zero gap, uniform contraction is impossible. This is target-wise
coverage at every structural model, not a claim of simultaneous
containment of the entire identified set at nonidentification.

The regimes are known, the means are known to be zero, the observations
are independent, and the dimension is fixed at two. Estimated regime
dates, unknown means under additional sampling conditions, non-Gaussian
identifying features and dependent observations are outside the theorem.
In particular a covariance argument does not identify rotations from
higher-order independent-component information when the covariance gap
is zero. The full joint identification modulus requested by the catalogue
and an optimal procedure for that broader experiment remain unresolved.

\section{Completion Estimate}
40\%

This is a heuristic estimate for the full catalogue target. The new
confidence construction removes both supplied-gap and known-whitening
restrictions from the earlier sharp Gaussian calculation, with a proved
uniform coverage and expected-diameter guarantee. The joint non-Gaussian,
dependent-data and unknown-regime problem remains outside the result.

\begin{thebibliography}{9}
\bibitem{lewis} D. J. Lewis.
\emph{Identification Based on Higher Moments in Macroeconometrics}.
Annual Review of Economics 17 (2025), 665--693.
\url{https://discovery.ucl.ac.uk/id/eprint/10208055/1/Lewis_annurev-economics-070124-051419.pdf}.
\bibitem{bhatia} R. Bhatia.
\emph{Matrix Analysis}.
Springer, 1997.
\end{thebibliography}
\end{document}
